\chapter{Introduction}

\section{Introduction}

\acp{LLM} were developed as a result of advancements in deep learning that allow them to read, understand and create large amounts of natural language. The development of these \acp{LLM} has been directly linked to an architectural advancement known as the Transformer. This architecture allows for the ability to model sequences based upon their relevance (attention) and it allowed for efficient training on larger than ever before data sets, in parallel \cite{vaswani2017attention}. With increases in both the amount of data available, computing power and model size, \acp{LLM} have achieved impressive results across a variety of tasks including question and answer, generating summaries, coding, and extracting information \cite{brown2020language,kaplan2020scaling}. These capabilities have provided several opportunities for utilizing them for various types of applications; however, with this comes significant technical barriers for deploying \acp{LLM} due to the issues associated with memory usage during processing, processing speed/latency for real-time operations, reliance upon other software packages, and ensuring reproducible results.

One technique used to improve both the utility and credibility of \ac{LLM} systems is called \acp{RAG}. The \ac{RAG} system takes advantage of a retrieval mechanism which pulls in relevant data from an outside body of documents, and then incorporates this retrieved data into the final input (i.e., model prompt) prior to generating an answer \cite{lewis2020rag,izacard2021leveraging}. Therefore, as opposed to answers being generated using only what is contained within the model's parameterized knowledge base, the answers are now supported by some form of evidence pulled directly from the outside world. In addition to supporting domain adaptation, \ac{RAG} can be done with no additional model training or full weight-fine tuning costs \cite{borgeaud2022retro}. 

At the same time as producing answers, \ac{RAG} systems aren't simply single-model systems. Document loading, chunking, embedding creation, indexing, retrieval (with optional reranking), prompt creation, answer creation, and log creation (i.e. logging results) all occur during the operation of a typical RAG system. Since each of these components can influence both the quality of the final answer created by the system, and the amount of time it takes to execute the entire \ac{RAG} system; therefore, in addition to determining if an answer was generated, a successful \ac{RAG} project will also need to determine how different retrieval strategies impact latency, evidence quality, answer accuracy, and reproducibility.

\ac{HPC} provides an interesting backdrop to this research. Because \ac{HPC} systems often consist of shared compute resources, job schedulers, containers for running code, large amounts of storage capable of supporting large-scale computation, and scalable hardware for running computational workloads \cite{kurtzer2017singularity,slurm2003,dongarra2020hpcsurvey}. Although this thesis limits its experimentation to the retrieval level \ac{RAG} techniques and how they compare to one another in terms of effectiveness and efficiency; the use of \ac{HPC} systems in providing frameworks for deployment requirements and scaling in terms of hardware and software to support larger scale workloads, frames our view of practical implementation as extending notebook-based reproducible pipeline workflows toward \ac{HPC} environments.

\section{Motivation}

The primary purpose behind creating this project is to address the increasing demand for language model applications which use private, domain-specific data or constantly evolving data. A single, self-contained language model will be able to generate coherent responses; however, the language model will likely lack immediate access to the specific documentation required to perform the desired function. Therefore, the model could potentially create responses which seem valid yet would be unsupported by existing documented proof. This is where \ac{RAG} can help alleviate some of these limitations. By allowing the \ac{RAG} system to search a previously established body of knowledge and retrieve relevant pieces of text, then pass these retrieved texts into the language model as contextual input prior to generating an answer.

Although the idea of \ac{RAG} seems relatively simple, the performance and speed of a full \ac{RAG} solution heavily depends upon how each piece of the solution is implemented. For example, building a very basic in memory retriever is easily accomplished and works well when working with smaller bodies of text. However, using a basic in memory retriever will not provide you with the same level of persistence or deployment options as a Vector Database. For instance, while a Vector Database like Milvus provides much better persistence options and a much more viable platform for working with large-scale retrieval requirements than an in-memory based approach; they come at an added expense through both set up costs and indexing costs. While reranking will allow you to reorder your retrieved items; reranking will add another model to your overall process and increase your overall retrieval time. Similarly, Hybrid Retrieval allows you to combine semantic and keyword-based retrieval methods; however, it will also add increased complexity to your overall retrieval method.

This leads us to why we created our experimental comparison-in-memory retrieval, vector-database retrieval, vector-database retrieval with optimized chunking, dense retrieval with reranking, hybrid retrieval and finally hybrid retrieval with reranking. We were also motivated by the need for a reproducible workflow due to many service providers, models, notebooks, configuration files, results logs, evaluation scripts and generated images make up a practical \ac{RAG} solutions.

\section{Problem Definition}

The general problem addressed by this Thesis is how to easilly and practically test different versions of a \ac{RAG} pipeline under a deployment workflow that is fully reproducible. The primary issue is not whether \ac{RAG} will improve general language-model applications as other studies have shown \cite{lewis2020rag,izacard2021leveraging}, but instead, what effect do different retrieval and ranking strategies have on speed, the amount and quality of retrieved evidence, and the final quality of answers and how easy it is to reproduce these results.

A \ac{RAG} pipeline contains several components whose behaviour must be coordinated. Documents must be loaded into memory and splitted into chunks; embedings must be generated; indexes or document stores must be created; candidate passages must be retrieved and finally, using the selected context, the language model must generate an answer. Weaknesses at any stage can reduce the usefulness of the final system. For example, while a fast retriever may return documents that are irrelevant to the user's query, these poorly choosen documents may lead to poor answers. Simultaneously, a re-ranking stage which may enhance the quality of the document choice may also result in significantly longer response times.

Therefore, the primary goal of this thesis is to study the implemented \ac{RAG} pipelines both as software engineering artifacts and experimental conditions. The goal of this study is to compare their output using exactly the same input documents, scenario queries, models and logging approach. This creates a controlled framework from with to evaluate the trade-off between performance and output quality.

The wider deployment environment of HPC has constraints such as job scheduling, container management, storage performance requirements and reproducibility requirements. In this thesis these constraints guide the design and documentation of the system. Large-scale multi-node benchmarking will be treated as future extensions.

\section{Research Objectives and Scope}

The main objective of this thesis is to develop \ac{RAG} pipeline prototypes and execute them in both a local setting and an \ac{HPC} environment in a reproducible manner. Since we do not aim to find the optimal settings for the model's weights, our focus will remain on the two components of the prototype that can be measured: retrieval methods at the document level and the ways of deploying the prototype as systems.

We attempt to use the same collections of documents, questions, models for languages and embeddings, as well as the same format for logging as much as possible throughout the dissertation. This enables us to measure the influence of different retrieval methods and deployment strategies independently. We have established a base line of behavior using the developed notebooks during the local phase. During the \ac{HPC} phase, the same files and scenarios were run so that we could compare results from local and HPC environments within a common evaluation framework.

To accomplish the goal outlined in this thesis, it will achieve the five following objectives:

\begin{enumerate}

    \item \textbf{Provide the theoretical and practical foundations for \acp{LLM} and \ac{RAG} Systems.}
    
    The thesis covers the technical foundation necessary for completing the task at hand which includes fine tuning, training from scratch, quantizing, pruning, adapting, etc., so that each of these background techniques can be differentiated from the various retrieval level techniques explored in this dissertation. Additionally, this section includes prompt engineering, \ac{RAG}, dense and sparse retrieval, hybrid retrieval, reranking, vector databases, and \ac{HPC} deployment.
    
    \item \textbf{Implement and assess the effectiveness of various \ac{RAG} Retrieval Methods.}
    
    The dissertation describes multiple methods for implementing retrieval: a baseline in-memory \ac{RAG} implementation; a Milvus-based dense-retrieval method utilizing a vector database; a modified version of the previously described method to optimize its performance; a method to use reranked answers with dense retrieval; a hybrid method to combine sparse and dense retrievals; and finally a hybrid method that combines sparse/dense retrievals and uses reranked answers. Each of the implementations is evaluated using standard retrieval metrics, such as hit rate, precision, recall, and F1, along with quality metrics for answers, such as manual correctness scores, ROUGE-L scores, BERTScores, and semantic similarity.
    
    \item \textbf{Determine the Performance Cost of the Implemented \ac{RAG} Pipelines.}
    
    This section provides the system-wide performance metrics for the \ac{RAG} pipelines that were implemented in the dissertation through notebook metrics and a log file: indexing time, retrieval time, and generation time. Using the above mentioned metrics will allow researchers to compare the costs associated with retrieving information in memory versus storing information in a database or disk. It also allows for a comparison of the benefits associated with reranking answers, combining sparse and dense retrievals, and optimized chunking.
    
    \item \textbf{Develop a Workflow to Compare Local vs. \ac{HPC} Execution.}
    
    Prior to executing the HPC portion of the dissertation research effort, local results are collected first. Once local results have been collected, notebooks for local execution are prepared for eventual execution on an HPC platform. The ultimate goal for the \ac{HPC} phase is to execute both local and \ac{HPC} versions of the dissertation's evaluation framework, allowing researchers to account for additional overhead factors such as scheduling delays, container startup times, storage access latency, and data transfer costs.
    
    \item \textbf{Create a Reproducible Evaluation Framework.}
    
    In addition to documenting the software tools utilized during execution, the dissertation documents all execution procedures, sample question sets used for testing, all fields used when logging results, manual evaluation criteria and processes used to manually evaluate results, scripts developed to ensure reproducibility of evaluations, and visualizations created during analysis.
\end{enumerate}

    This research limits its study to inference time \ac{LLM} and \ac{RAG} loads. This will enable an experiment to focus on retrieving quality answers as well as system performance. The scope does NOT include full model training or large scale distributed training, model fine tuning, model compression (quantization) or pruning, model distillation, and adapter based model optimization. These concepts have been referenced where necessary for a comparison in concept.

\section{Research Questions and Hypotheses}

The research questions are designed to remain within the experimental scope of the thesis. They focus on two measurable areas: the accuracy and answer quality of retrieval-level techniques, and the performance impact of system-level strategies. The local experiments provide the first set of measurements, while the later \ac{HPC} experiments repeat the same workflow to evaluate portability and deployment behaviour.

\begin{itemize}

\item {\bfseries RQ1:} 
\textbf{How do different retrieval techniques affect the retrieval-accuracy and answer-quality of the implemented \ac{RAG} pipelines?}
{\bf H1: } 
Hybrid-retrieval will reach a better retrieval-coverage than dense-retrieval due to lexical-matching combined with semantic-similarity; Reranking will make sure that retrieved passages are ordered well and thus increase top-k precision, answer-groundedness, and answer-correctness. It is assumed that the combination of hybrid-retrieval and reranking will lead to the best retrieval-quality-results particularly for those queries which need an understanding of their semantic interpretation as well as the correct wording. Another expected noticable improvement is the tuned Vector database implementation.

\item {\bfseries RQ2: }
\textbf{What kinds of trade-off regarding the performance are introduced through hybrid-retrieval and reranking?}
{\bfseries H2: } 
The advantages at the level of retrieval are not cost-free in terms of execution-time. As a result, it is assumed that the retrieval-latency of the two basic (vector-database) pipelines will be lower than the one introduced by reranking and by hybrid-retrieval-plus-reranking. Consequently, the best-performing pipeline concerning answer-quality might not be the fastest pipeline concerning end-to-end latency.

\item {\bfseries RQ3: }
\textbf{How do system-level strategies like caching, batching, and parallelism influence the performance of \ac{RAG} pipelines?}
{\bfseries H3: } 
Caching is assumed to lower repeated query retrieval-time as well as repeated index loading-overhead without changing the retrieval accuracy. Batching is assumed to improve the throughput if several queries are processed simultaneously. However, there is also a risk that batching leads to increased latencies for individual queries under low workload conditions. Parallelism is expected to reduce the indexing time and the wallclock time of batch evaluations more clearly than it reduces the single query generation latency since the generation still depends on the runtimes of the language model.

\item {\bfseries RQ4: }
\textbf{How behaves the same \ac{RAG} workflow when moved from local-deployment to \ac{HPC} deployment?}
{\bfseries H4: } 
If the same files, scenarios, models, and log-format are employed, then retrieval accuracy and answer quality are expected to stay relatively similar in local and \ac{HPC} execution. Differences are expected to show up mainly in performance and operationally. On \ac{HPC}, it is expected to improve throughput and batch-execution when parallelism and batching are employed. However, single-query latency might not always benefit since scheduling-delays, container-startup-times, shared-storage-access times, and data-transfer overheads could compensate for the additional computing-resources available.

\end{itemize}

These hypotheses are deliberately comparative and quantifiable. Every one of them can be assessed utilizing the established evaluation framework: retrieval-quality metrics, answer-quality rubrics, timing fields, throughput measurements, resource measurements, and environment labels for local and \ac{HPC} execution.

\section{Research Methodology Overview}

\begin{figure}[htbp]
    \centering
    \includegraphics[width=\textwidth]{images/fig01_overall_thesis_workflow.png}
    \caption{Overall thesis workflow showing the local deployment stage, the later HPC deployment stage, and the final comparison between \ac{RAG} variants and execution environments.}
    \label{fig:overall-thesis-workflow}
\end{figure}

The methodology includes a literature study, a software implementation, controlled experiments, and an analysis of the qualitative output from these studies.

The 1st Step develpos the theoretical and technological basis for the studies. It reviews \acp{LLM}, transformer architectures, how prompts are generated, \ac{RAG} architectures, retrieval technologies, reranking, hybrid retrieval and \ac{HPC} as a tehcnology for deploying systems. Step one provides the vocabulary and context necessary to understand both the implementation (Step 2) and the results (Step 2 and 3).

Step 2 represents the implementation of the research pipeline. The task-related aspect of this project was organized into a series of six Jupyter notebooks that implement a basic \ac{RAG} pipeline, a vector database pipeline, a fine-tuned vector database pipeline, a reranking pipeline, a hybrid retrieval pipeline and finally, a hybrid retrieval with reranking pipeline. Following the same general format as one another, the content of each notebook includes; loading documents, preparing chunks, creating or connecting to the appropriate store (vector database), running query scenarios, generating answers and saving results.

The 3rd step represents the evaluation process for each of the pipelines. To evaluate each pipeline, timing information is collected. The timing information collected from each of the pipelines is based on seconds. Indexing time, retrieval time and generation time are all included in this measurement. Additionally to collecting qualitative measurements of how well each pipeline performed, the source documents retrieved and the generated answers were also manually scored and analyzed utilizing deterministic metrics. The manual scores represent if the generated answer is correct or incorrect as determined by the evidence within the document corpus. Deterministic metrics utilized include retrieval hit, retrieval precision, retrieval recall, retrieval F1 score, ROUGE-L, BERTScore, semantic similarity and length of answer. Collecting both qualitative and quantitative measurements allows for an assessment of the trade-offs referenced in Chapter 6.

The three primary research objectives developed into this thesis were:


\section{Expected Contributions}

\begin{enumerate}
    \item {\textbf{Development of a local, reproducible framework for evaluating \ac{RAG} pipelines, which will serve as a basis for later high-performance computing (\ac{HPC}) evaluations}} 
    The authors provide a defined process for running different \ac{RAG} implementations using the same set of input documents, question scenarios, log parameters and measurement methods. While all of the empirical data used in the current study has been gathered at the local level; the approach described here could be applied to the HPC platform in future studies.
    
    \item {\textbf{Empirical comparison of \ac{RAG} techniques at the document level}} 
    In addition to comparing basic in-memory search (retrieval), Milvus-based vector database search, optimized vector database search, dense search with reranking, hybrid sparse-dense search, and hybrid search with reranking; we compare them based on the metrics that have been built into our project: number of hits found by the retrieval mechanism, number of relevant items returned by the retrieval algorithm, number of correct answers provided by the retrieval algorithm, retrieval F1, manually scored answers, ROUGE-L scores, BERTScore values, similarity values between semantically equivalent concepts, and lengths of retrieved answers.
    
    \item {\textbf{Practical assessment of the trade-offs between quality and efficiency in executing \ac{RAG} locally.}} 
    We measure how the choice of retrieval impacts the times required to index documents, retrieve information from those indexes, generate responses based on the retrieved information and produce retrievable evidence.
\end{enumerate}

Future directions include investigating larger-scale system-wide alternatives (e.g., caching mechanisms, batch processing capabilities, ability to execute multiple processes concurrently, measuring throughput of individual components, memory usage monitoring) as well as comparing execution at both the local and HPC levels.

\section{Structure of the Thesis}

The rest of the thesis will be structured as follows:

\textbf{Chapter 2 }will discuss the theory and history of the background. The chapter covers:
\acp{LLM}, the architecture of Transformer, training and adaptation strategies for \acp{LLM}, prompt engineering, the \ac{RAG} architecture, retrieval methods, and evaluation metrics for \ac{HPC}.

\textbf{Chapter 3} will review other works about:
deployment of \acp{LLM} (including their optimizations during inference), systems based on \ac{RAG} architectures, techniques to retrieve information, and benchmarking methods. Chapter 3 identifies what motivated the implementation of the strategy for the evaluation presented in this thesis.

\textbf{Chapter 4} will present a description of methodology. This includes the research framework, input data, system architecture, implemented pipelines based on \ac{RAG} architectures, experimental parameters, performance measures, answer-quality measures, and reproducibility procedures.

\textbf{Chapter 5} provides an overview of how experiments were set up and executed. The chapter provides details on the software platform, models and services, input resources, the order in which notebooks are executed, experiment scenarios, and logging procedure for results.

\textbf{Chapter 6} presents and examines results. The chapter compares the various pipelines implemented with regard to time-based metrics and output-quality metrics; it also analyzes the trade-off between speed, retrieval quality, answer quality, and reproducibility.

\textbf{Chapter 7} will summarize the conclusions from this thesis. The chapter will outline the main contributions of this thesis; it will provide examples of its limitations; and it will propose potential future directions for this area of study, especially large scale deployments of HPC systems, additional datasets, and more advanced evaluation methodologies.

\section{Summary}

This Chapter discussed the reason for this study, the problem, the objective of this study, the Research Questions, Methodology, Potential Contributions, and Structure of the Thesis. In order to further narrow down the scope of this study to be only on the \ac{RAG} pipelines that were implemented, as well as their trade-off's in terms of what is retrievable at each level (retrieval), the quality of outputs produced by these pipelines, and Reproducibility. The next Chapter will cover all of the Theory behind Language Models, Retrieval models and \ac{HPC} systems that are necessary for a thorough understanding of the concepts used within this project.
